\section{Methodology}

\subsection{Experimental Design}

\begin{table}[h]
\centering
\begin{tabular}{llll}
\toprule
ID & Hypothesis & Gate & Purpose \\
\midrule
H-E1 & Benchmark Independence & MUST\_WORK & Validate experimental design \\
H-M1 & RLHF Reward Smoothing & MUST\_WORK & Confirm RLHF mechanistic claim \\
H-M2 & DPO Boundary Sharpness & SHOULD\_WORK & Confirm DPO mechanistic claim \\
H-M3 & Different Attractors & SHOULD\_WORK & Test landscape$\rightarrow$attractor theory \\
H-M4 & Differential Profiles & SHOULD\_WORK & Test attractor$\rightarrow$benchmark claim \\
\bottomrule
\end{tabular}
\caption{Experimental design with five hypotheses.}
\label{tab:hypotheses}
\end{table}

\subsection{Setup}

\textbf{Model:} Llama-2-7B \citep{touvron2023llama} with LoRA \citep{hu2022lora} ($r=16$, $\alpha=32$)

\textbf{Data:} Anthropic HH-RLHF (15k training subset)

\textbf{Evaluation:} TruthfulQA MC1 (817), HHH-helpful (1000), HHH-harmless (1000)

\subsection{Metrics}

\begin{itemize}
    \item \textbf{H-M1:} Reward output range (smoothness)
    \item \textbf{H-M2:} Sharpness ratio = DPO margin variance / RLHF margin variance
    \item \textbf{H-M3:} Clustering gap = within-method similarity $-$ cross-method similarity
    \item \textbf{H-M4:} Differential profile = max$|d| > 0.3$ AND min$|d| < 0.15$
\end{itemize}
